\documentclass{article}%
\usepackage[T1]{fontenc}%
\usepackage[utf8]{inputenc}%
\usepackage{lmodern}%
\usepackage{textcomp}%
\usepackage{lastpage}%
\usepackage{geometry}%
\geometry{margin=1in,headheight=14pt,headsep=25pt}%
\usepackage{hyperref}%
\usepackage{enumitem}%
\usepackage{fancyhdr}%
\usepackage{xcolor}%
\usepackage{url}%
\usepackage{breakurl}%
%

            \hypersetup{
                colorlinks=true,
                linkcolor=blue,
                filecolor=magenta,
                urlcolor=blue
            }
            \pagestyle{fancy}
            \fancyhf{}
            \rhead{Generated on \today}
            \cfoot{\thepage}
            
            % Configure enumeration settings
            \setlist[enumerate,1]{label=\arabic*.}
            \setlist[enumerate,2]{label=\alph*.}
            \setlist[enumerate,3]{label=\Alph*.}
            \setlist[enumerate]{itemsep=0.5em}
            
            % Define a command for red text
            \newcommand{\incorrect}[1]{\textcolor{red}{#1}}
        %
\lhead{Complementary Slackness}%
\title{Practice Questions for Complementary Slackness}%
\author{Generated by Scribe.AI}%
\date{\today}%
%
\begin{document}%
\normalsize%
\maketitle%
\section{Questions}%
\label{sec:Questions}%
\begin{enumerate}%
\item In the context of complementary slackness, what is the relationship between the primal and dual variables if a constraint in the primal problem is non-binding (i.e., the inequality holds strictly)?%
\begin{enumerate}%
\item The corresponding dual variable must be zero.%
\item The corresponding dual variable must be positive.%
\item The corresponding dual variable can be either zero or positive.%
\item The corresponding primal variable must be zero.%
\item There is no relationship between the primal and dual variables in this case.%
\end{enumerate}%
\vspace{1em}%
\item Explain the concept of complementary slackness and its significance in verifying the optimality of solutions in linear programming.%
\begin{enumerate}%
\item Complementary slackness states that if a primal variable is positive, the corresponding dual constraint must be binding, and vice versa.%
\item Complementary slackness is a condition that must hold only for infeasible linear programs.%
\item Complementary slackness is an alternative method for solving linear programs, replacing the simplex method.%
\item Complementary slackness is only relevant for dual problems, not primal problems.%
\item Complementary slackness is unrelated to the optimality of solutions in linear programming.%
\end{enumerate}%
\vspace{1em}%
\item Suppose you have solved a primal linear programming problem and found an optimal solution. How can complementary slackness be used to verify the optimality of the corresponding dual solution?%
\begin{enumerate}%
\item Complementary slackness is irrelevant for verifying the dual solution's optimality.%
\item Check if the primal objective function value equals the dual objective function value.%
\item Verify that the complementary slackness conditions are satisfied for both the primal and dual solutions.%
\item Solve the dual problem independently and compare the objective function values.%
\item There is no way to verify the dual solution's optimality using only the primal solution.%
\end{enumerate}%
\vspace{1em}%
\item What is the practical implication of complementary slackness in sensitivity analysis of linear programming problems?%
\begin{enumerate}%
\item Complementary slackness has no role in sensitivity analysis.%
\item It helps determine how changes in the objective function coefficients affect the optimal solution.%
\item It helps determine which constraints are binding and which are non-binding at the optimal solution, aiding in understanding the impact of changes to resource availability or constraint parameters.%
\item It simplifies the calculation of the shadow prices but provides no information about the binding constraints.%
\item It is only useful for identifying unbounded solutions.%
\end{enumerate}%
\vspace{1em}%
\item Explain the concept of complementary slackness in linear programming and its significance in verifying the optimality of solutions.%
\begin{enumerate}%
\item Complementary slackness states that if a constraint in the primal problem is binding, then the corresponding dual variable must be zero.  If a constraint is non-binding, the corresponding dual variable can be any non-negative value.%
\item Complementary slackness is a theorem that guarantees the existence of an optimal solution for both the primal and dual problems, provided that both are feasible.%
\item Complementary slackness is a condition that must be satisfied by any feasible solution to the primal and dual problems, regardless of whether they are optimal.%
\item Complementary slackness states that at optimality, the product of each primal variable and its corresponding dual slack variable is zero, and the product of each dual variable and its corresponding primal slack variable is zero.%
\item Complementary slackness is a method for converting a primal linear programming problem into its dual form.%
\end{enumerate}%
\vspace{1em}%
\item In the context of complementary slackness, if the $i$-th constraint in the primal problem is non-binding (i.e., the inequality holds strictly), what can we say about the corresponding dual variable $y_i$?%
\begin{enumerate}%
\item $y_i$ must be equal to zero.%
\item $y_i$ must be greater than zero.%
\item $y_i$ can be any non-negative value.%
\item $y_i$ must be less than zero.%
\item $y_i$ can be any value.%
\end{enumerate}%
\vspace{1em}%
\item Consider a primal linear programming problem and its dual.  If the primal problem has an optimal solution with an objective function value of $Z^*$, and the dual problem has an optimal solution with an objective function value of $W^*$, what does strong duality tell us about the relationship between $Z^*$ and $W^*$?%
\begin{enumerate}%
\item $Z^* \ge W^*%
\item $Z^* \le W^*%
\item $Z^* = W^*%
\item $Z^* > W^*%
\item $Z^* < W^*%
\end{enumerate}%
\vspace{1em}%
\item Explain how complementary slackness can be used to verify the optimality of a solution to a linear programming problem and its dual.%
\begin{enumerate}%
\item Complementary slackness is not relevant to verifying optimality.%
\item If a primal constraint is binding, the corresponding dual variable must be zero.%
\item If a primal variable is positive, the corresponding dual constraint must be binding.%
\item If a primal constraint is non-binding, the corresponding dual variable must be positive.%
\item Complementary slackness states that either a primal constraint is binding or the corresponding dual variable is zero, but not both.%
\end{enumerate}%
\vspace{1em}%
\item Suppose you have the primal problem: Maximize $Z = 2x_1 + 3x_2$ subject to $x_1 + x_2 \le 5$, $x_1 \le 3$, $x_2 \le 4$, $x_1, x_2 \ge 0$.  An optimal solution is found to be $x_1 = 2, x_2 = 3$, with $Z = 13$.  The corresponding dual problem has optimal solution $y_1 = 1, y_2 = 1, y_3 = 0$. Verify the optimality using complementary slackness.%
\begin{enumerate}%
\item Complementary slackness is not satisfied.%
\item Complementary slackness is satisfied for $x_1$ and $y_2$ but not for $x_2$ and $y_3$.%
\item Complementary slackness is satisfied for $x_2$ and $y_3$ but not for $x_1$ and $y_2$.%
\item Complementary slackness is satisfied.%
\item The given primal solution is not optimal.%
\end{enumerate}%
\vspace{1em}%
\item Consider the primal problem: Maximize $Z = 3x_1 + 2x_2$ subject to $-x_1 + 3x_2 \le 12$, $x_1 + x_2 \le 8$, $2x_1 - x_2 \le 10$, $x_1, x_2 \ge 0$.  Suppose an optimal solution is $x_1 = 4, x_2 = 4$ with $Z = 20$. The corresponding dual problem has an optimal solution $y_1 = 1, y_2 = 1, y_3 = 0$. Verify the optimality using complementary slackness. What is the value of $x_1 z_1$?%
\begin{enumerate}%
\item 0%
\item 4%
\item 12%
\item 20%
\item Cannot be determined%
\end{enumerate}%
\vspace{1em}%
\end{enumerate}

%
\newpage%
\section{Answers}%
\label{sec:Answers}%
\begin{enumerate}%
\item In the context of complementary slackness, what is the relationship between the primal and dual variables if a constraint in the primal problem is non-binding (i.e., the inequality holds strictly)?%
\begin{enumerate}%
\item If a primal constraint is non-binding, the corresponding dual variable must be zero. This is a direct consequence of the complementary slackness theorem.%
\item \incorrect{A positive dual variable implies that the corresponding primal constraint is binding (holds with equality).}%
\item \incorrect{The dual variable cannot be positive if the primal constraint is non-binding.}%
\item \incorrect{The primal variable's value is independent of whether the corresponding dual constraint is binding or not.}%
\item \incorrect{Complementary slackness establishes a direct relationship between primal constraints and dual variables.}%
\end{enumerate}%
\vspace{1em}%
\item Explain the concept of complementary slackness and its significance in verifying the optimality of solutions in linear programming.%
\begin{enumerate}%
\item Complementary slackness provides necessary and sufficient conditions for primal and dual solutions to be optimal.  It links the primal variables and dual constraints (and vice versa).%
\item \incorrect{Complementary slackness applies to both feasible and infeasible linear programs, but it's particularly useful for verifying optimality in feasible cases.}%
\item \incorrect{Complementary slackness is a theorem, not a solution method. It's used to check the optimality of solutions found by other methods like the simplex method.}%
\item \incorrect{Complementary slackness applies to both primal and dual problems, establishing a relationship between them.}%
\item \incorrect{Complementary slackness is fundamental to verifying the optimality of solutions in linear programming.}%
\end{enumerate}%
\vspace{1em}%
\item Suppose you have solved a primal linear programming problem and found an optimal solution. How can complementary slackness be used to verify the optimality of the corresponding dual solution?%
\begin{enumerate}%
\item \incorrect{Complementary slackness is crucial for verifying the optimality of both primal and dual solutions.}%
\item \incorrect{While equality of objective function values is a consequence of strong duality, it doesn't directly verify complementary slackness conditions.}%
\item Complementary slackness provides a direct link between primal and dual optimal solutions.  Checking the conditions ensures both are optimal.%
\item \incorrect{Solving the dual independently is redundant; complementary slackness offers a more efficient verification method.}%
\item \incorrect{Complementary slackness provides a mechanism to verify the dual solution's optimality using the primal solution.}%
\end{enumerate}%
\vspace{1em}%
\item What is the practical implication of complementary slackness in sensitivity analysis of linear programming problems?%
\begin{enumerate}%
\item \incorrect{Complementary slackness is a key tool in sensitivity analysis, providing insights into the impact of changes in constraints.}%
\item \incorrect{While changes in objective function coefficients affect the optimal solution, complementary slackness primarily focuses on the constraints.}%
\item Understanding which constraints are binding (and thus the corresponding dual variables) is crucial for sensitivity analysis. Complementary slackness directly helps identify these.%
\item \incorrect{Complementary slackness helps identify both shadow prices (dual variables) and binding constraints.}%
\item \incorrect{Complementary slackness is a general tool for optimality verification and sensitivity analysis, not limited to unbounded solutions.}%
\end{enumerate}%
\vspace{1em}%
\item Explain the concept of complementary slackness in linear programming and its significance in verifying the optimality of solutions.%
\begin{enumerate}%
\item \incorrect{This option partially describes complementary slackness but is incomplete and inaccurate. It doesn't cover the condition for dual variables and primal slack variables.}%
\item \incorrect{While strong duality guarantees the equality of optimal objective function values if both primal and dual problems are feasible, complementary slackness is a separate condition related to the relationship between primal and dual variables at optimality.}%
\item \incorrect{Complementary slackness is not a condition for any feasible solution; it's a condition that holds only at optimality.}%
\item Complementary slackness is a crucial condition that must hold at optimality for both the primal and dual problems.  It states that the product of each primal variable and its corresponding dual slack variable is zero, and similarly for the dual variables and primal slack variables. This condition provides a powerful tool for verifying the optimality of a solution pair for the primal and dual problems.%
\item \incorrect{Complementary slackness is a concept related to the relationship between primal and dual solutions, not a method for converting between them.}%
\end{enumerate}%
\vspace{1em}%
\item In the context of complementary slackness, if the $i$-th constraint in the primal problem is non-binding (i.e., the inequality holds strictly), what can we say about the corresponding dual variable $y_i$?%
\begin{enumerate}%
\item Complementary slackness states that if a primal constraint is strictly satisfied (non-binding), then the corresponding dual variable must be zero.  This is because there is no "slack" to be priced in the dual problem.%
\item \incorrect{If $y_i > 0$, it implies that the $i$-th primal constraint is binding (active), not non-binding.}%
\item \incorrect{While $y_i$ must be non-negative, it cannot be any non-negative value if the corresponding primal constraint is non-binding. It must be zero.}%
\item \incorrect{Dual variables are always non-negative.}%
\item \incorrect{The value of $y_i$ is restricted by complementary slackness; it cannot be any value.}%
\end{enumerate}%
\vspace{1em}%
\item Consider a primal linear programming problem and its dual.  If the primal problem has an optimal solution with an objective function value of $Z^*$, and the dual problem has an optimal solution with an objective function value of $W^*$, what does strong duality tell us about the relationship between $Z^*$ and $W^*$?%
\begin{enumerate}%
\item \incorrect{This describes weak duality, which states that the primal objective function value is always greater than or equal to the dual objective function value for any feasible solutions.}%
\item \incorrect{This is the reverse of weak duality.}%
\item Strong duality states that if both the primal and dual problems have optimal solutions, then their optimal objective function values are equal.  That is, $Z^* = W^*$.%
\item \incorrect{This is incorrect; strong duality states that the optimal values are equal.}%
\item \incorrect{This is incorrect; strong duality states that the optimal values are equal.}%
\end{enumerate}%
\vspace{1em}%
\item Explain how complementary slackness can be used to verify the optimality of a solution to a linear programming problem and its dual.%
\begin{enumerate}%
\item \incorrect{Complementary slackness is a fundamental concept in linear programming for verifying optimality.}%
\item \incorrect{This is incorrect; if a primal constraint is binding, the corresponding dual variable can be positive.}%
\item \incorrect{This is the correct relationship, but it's only one part of complementary slackness. The complete condition involves both this and the condition for non-binding constraints.}%
\item \incorrect{If a primal constraint is non-binding, the corresponding dual variable must be zero.}%
\item Complementary slackness provides necessary and sufficient conditions for optimality.  It states that for each pair of a primal constraint and its corresponding dual variable, either the constraint is binding (equality holds) or the dual variable is zero, but not both.  The same applies to primal variables and dual constraints.%
\end{enumerate}%
\vspace{1em}%
\item Suppose you have the primal problem: Maximize $Z = 2x_1 + 3x_2$ subject to $x_1 + x_2 \le 5$, $x_1 \le 3$, $x_2 \le 4$, $x_1, x_2 \ge 0$.  An optimal solution is found to be $x_1 = 2, x_2 = 3$, with $Z = 13$.  The corresponding dual problem has optimal solution $y_1 = 1, y_2 = 1, y_3 = 0$. Verify the optimality using complementary slackness.%
\begin{enumerate}%
\item \incorrect{The provided solution is not optimal, but complementary slackness can still be checked.}%
\item \incorrect{Incorrect analysis of complementary slackness.}%
\item \incorrect{Incorrect analysis of complementary slackness.}%
\item \incorrect{}%
\item \incorrect{The given primal solution is not optimal because complementary slackness is not satisfied.}%
\end{enumerate}%
\vspace{1em}%
\item Consider the primal problem: Maximize $Z = 3x_1 + 2x_2$ subject to $-x_1 + 3x_2 \le 12$, $x_1 + x_2 \le 8$, $2x_1 - x_2 \le 10$, $x_1, x_2 \ge 0$.  Suppose an optimal solution is $x_1 = 4, x_2 = 4$ with $Z = 20$. The corresponding dual problem has an optimal solution $y_1 = 1, y_2 = 1, y_3 = 0$. Verify the optimality using complementary slackness. What is the value of $x_1 z_1$?%
\begin{enumerate}%
\item \incorrect{}%
\item \incorrect{This is incorrect.  Complementary slackness requires that if a primal variable is positive, the corresponding dual constraint is binding, and vice versa.  The calculation is incorrect.}%
\item \incorrect{This is incorrect. The calculation is incorrect. Complementary slackness conditions are not satisfied.}%
\item \incorrect{This is incorrect. The calculation is incorrect. Complementary slackness conditions are not satisfied.}%
\item \incorrect{This is incorrect. Complementary slackness provides a way to verify the optimality. The calculation can be performed.}%
\end{enumerate}%
\vspace{1em}%
\end{enumerate}

%
\end{document}